\chapter{Ontology of Harmonic Information}

\epigraph{\emph{A single oscillator can oscillate. Two oscillators can enter into relation. From relation comes interference; from interference, pattern; from pattern, the possibility of differential uptake.}}{}


If the previous chapters established why a coordinated field may now be warranted, the present chapter addresses a prior question: what, exactly, is the object that such a field would study? Harmonic Information Theory cannot remain a useful category of coordination unless it specifies, at least provisionally, what it means by \texttt{information}, why harmonic relations might matter ontologically, and under what conditions an interval can be treated as more than a descriptive convenience. The immediate task is therefore conceptual precision: to define the vocabulary without which later claims about consonance, efficiency, or biological significance would remain equivocal.

The central proposal is deliberately minimal. HIT begins from the possibility that information is not exhausted by symbolic encoding, probabilistic distribution, or the properties of isolated substances. Under some conditions, information may also reside in structured relations among oscillatory, resonant, and wave-bearing processes. More specifically, an interval - understood first as a ratio relation between frequencies or modes and only secondarily as a musical interval in the narrow sense - may function as a minimal carrier of organization insofar as it generates recurrent, scale-robust, and differentially interpretable patterns. This is not yet a completed theory. It is an ontological proposal about what kinds of patterns deserve to count as primary in the first place.


\section{Information as pattern, not only as quantity}


Any ontology of harmonic information must begin by resisting a common simplification: the assumption that \texttt{information} already has one settled meaning and that all further uses are metaphorical deviations from it. In fact, the modern concept of information is already stratified. Shannon's theory remains indispensable because it formalized information as reduction of uncertainty over a space of possible messages \parencite{shannon_1948}. That achievement was decisive. It made information measurable, transmissible, and engineerable. But Shannon's abstraction was also intentionally restrictive: it bracketed semantics, embodiment, and function in order to isolate the formal constraints of communication. That move gave us a rigorous measure of informational quantity. It did not, by itself, settle what kinds of patterns deserve ontological priority in natural systems.

Once one leaves the engineering frame, a second layer becomes unavoidable. Information is not only what reduces uncertainty in a channel; it is also what exhibits structure rather than arbitrariness. In algorithmic and compressibility-oriented perspectives, a pattern is informative not because it is maximally random, but because it admits a shorter and more lawful description than an unstructured aggregate would. Put differently, structure matters because it can be compactly represented, predicted, or regenerated. This does not negate Shannon; it supplements it. The point is that not all informative patterns are equally primitive. Some carry lawlike regularities that allow systems to coordinate with lower descriptive or predictive cost. Dubnov's work on information and musical expectation points in that direction by tying salience to structured predictability rather than to raw uncertainty alone \parencite{dubnov_2004}.

Bateson adds a third layer that is essential for HIT. His well-known formulation of information as "a difference which makes a difference" relocates the question from mere formal distribution to operative consequence \parencite{bateson_1972}. In this register, a pattern counts as informational only if it can alter the state, expectation, or behavior of some receiving system. The advantage of this formulation is that it prevents a purely reified ontology of form. A pattern is not informative simply because an observer can describe it elegantly. It is informative when it is capable of differential uptake, constraint, or effect. This point will matter later when HIT moves from formal interval structure to biological and computational interpretation. At the ontological level, however, the lesson is already clear: information is relational not only because it consists of differences, but because it exists through the possibility of system-level response.

The fourth layer is oscillatory and dynamical. In many natural systems, information does not travel only through symbolic tokens or static states; it is enacted through coupling, recurrence, synchronization, phase structure, propagation, and mode organization. Kuramoto's framework made one part of this intuition formally compelling by showing that some relations among oscillators are dynamically privileged: simple ratios can stabilize coordination regimes that would not emerge under arbitrary mismatch \parencite{kuramoto_1984,strogatz_2000}. In neuroscience, cross-frequency coupling extends this point into living systems: oscillatory relations are not incidental decorations on top of information processing, but often part of the very mechanism by which integration occurs across scales \parencite{canolty_2010}. In physics, the modern vocabulary of fields, waves, and excitations preserves the same structural lesson in a different idiom: what matters is often not the isolated bearer but the patterned organization of modes, their couplings, and their interference.

This does not license the crude slogan that everything is simply "wave" in one undifferentiated sense. Different theories retain different ontological commitments, and no single model should be inflated into the whole of what its own discourse constructs as reality. But it does make one claim newly reasonable: across a wide range of domains, some of the most powerful descriptions of organization are descriptions of patterned relation in time and space. Information, in this register, is associated not only with uncertainty reduction or compressibility, but with the emergence of stable, recurrent, and interpretable organizations among coupled processes.

These four layers should not be collapsed into one another. Shannon is not Bateson; compressibility is not biological meaning; dynamic stability is not yet semantic consequence. But they do converge on a common intuition: information is not exhausted by isolated substances or by raw magnitudes taken one at a time. It is tied to patterned difference, lawful relation, and the possibility of repeatable uptake. For the purposes of HIT, this convergence is enough to justify a working definition. We will treat information, in the broad ontological sense relevant to this project, as a structured relation or detectable pattern in a signal or system that, for a given observer, system, or model, can reduce uncertainty, increase compressibility, and produce a functional difference through recurrent interpretability. Under suitable physical conditions, such patterned relation may also participate in local reductions of effective entropy, but that consequence should not be confused with uncertainty reduction in the narrower Shannon sense. What becomes legible as information is therefore not a transparent window onto the Real, but a structured trace through which the Real imposes itself as limit, resistance, or failure of closure within a model. Harmonic information is then a restricted case of this broader category: information carried, constrained, or stabilized by harmonic relations among oscillatory or mode-bearing processes.


\section{Why relation can be ontologically primary}


This definition becomes intelligible only if relation is allowed genuine ontological weight. Much of modern scientific common sense still leans, often implicitly, on a substance-first picture: things are primary, and relations are later properties among already constituted units. HIT requires a different emphasis. The claim is not that substances dissolve into pure relation, but that some of the most durable regularities in the domains under discussion appear first as patterned coupling, interference, and co-constraint. In such cases, relation is not decorative. It is part of what makes the phenomenon legible at all.\footnote{For a recent engineering instance where information appears as physically realized geometry rather than as a detachable abstract layer, see Jelin{\v{c}}i{\v{c}} et al.\ (2025). Their probabilistic hardware formalism treats the accessibility of states as a consequence of energy-landscape geometry and relative transition structure, not as a property of isolated states considered one by one.}

Heidegger's account of being-in-the-world begins to unsettle the inherited picture at exactly this point. Entities do not first present themselves as inert objects and only afterward acquire significance through external linkage. They emerge within a prior field of worldhood, practice, discourse, and being-with, where relevance and articulation are already in play \parencite{heidegger_1927}. Bohm approaches a related displacement from another direction when he argues that what appears as separation may be the unfolding of a deeper implicate order rather than the final truth of the situation \parencite{bohm_1980}. Simondon and Barad push the point into explicitly relational vocabularies. Individuation, transduction, and intra-action all name processes in which units do not simply precede their relations, but take shape through them \parencite{simondon_1958,barad_2007}. Taken together, these traditions do not erase the world of things. They make it harder to assume that things are self-sufficient first terms and that relation arrives only afterward as a secondary predicate.

Contemporary consciousness-first ontologies push this displacement further. Hoffman's conscious realism treats spacetime as an interface generated by conscious agents, while Kastrup's analytic idealism treats matter as the extrinsic appearance of mentation within a wider field of consciousness \parencite{hoffman_2008,kastrup_2017a,kastrup_2017b}. Whether one accepts these inversions or not, they matter here for a narrower reason. They show that the priority of relation, appearance, and experiential articulation over self-sufficient substance is not confined to continental philosophy or systems biology. It has become an active site of formal and ontological experimentation.

Maturana and Varela provide a biological formulation of the same displacement. In the autopoietic account, a living system is not defined first by the inventory of its material parts, but by the organization of relations through which those parts are continuously produced and maintained \parencite{maturana_varela_1980,varela_maturana_uribe_1974}. This is why their distinction between organization and structure matters so much here. Organization names the invariant relational pattern that defines the system as the kind of unity it is; structure names the concrete material realization through which that organization is instantiated at a given moment \parencite{maturana_1970,varela_1979}. Read cautiously, HIT can borrow a narrow analogy from this distinction. A ratio may function like an organizational invariant, while the particular frequencies, amplitudes, and phases that realize it belong to the order of structure. The point is not that harmonic ratios are living systems. It is that Maturana and Varela supply one of the clearest biological precedents for treating relational pattern as prior to the isolated identity of components.

Operational closure sharpens the same lesson. An autopoietic system is materially open yet organizationally closed in the sense that its operations recursively generate the very network that makes those operations possible. This circularity is not self-identity in the static sense, but productive recurrence. HIT can draw a carefully limited comparison from that form of closure to the way some interference patterns or standing-wave organizations preserve a recognizable relational form under changing energetic conditions. The comparison must remain explicitly limited. A standing wave does not regenerate its own material components and is therefore not autopoietic in the strict sense. What it does share with autopoietic description is a focus on relational invariance under structural variation. That is enough to make Maturana and Varela newly relevant to the ontological wager of this chapter without collapsing biological autonomy into harmonic pattern.

The same problem also becomes visible from outside biological theory proper, which matters here because it shows that recursive organization is not a local intuition of one school alone.

An instructive convergence appears here from a different tradition altogether. In \emph{La tecnica digitalizada}, Fern\'andez M\'endez describes a "relational closure" through which a system of know-how achieves enough internal consistency to sustain itself without the continuous intervention of a subject \parencite{fernandez_2021b}. In the manuscript on digital language, he arrives at a closely related problem from yet another route, connecting recursive formal systems, Godelian incompleteness, and operational closure as conditions under which a system can function without ever becoming complete in the strong sense \parencite{fernandez_sf_digital}. HIT does not need to collapse these trajectories into one doctrine. It only needs to register a structurally significant fact: biology, psychoanalytic theories of work and technique, and the epistemology of computation all converge on the intuition that organization is achieved through recursive relation before it is secured by isolated substantial units.

That shift matters for HIT because the framework is interested in phenomena whose explanatory core often lies in organization among terms rather than in the isolated identity of those terms. A ratio does not become significant merely because two frequencies can be listed side by side; it becomes significant when their relation constrains pattern, recurrence, interference, and uptake. This is why the issue is ontological before it is methodological. If relation were only derivative, harmonic information would always risk appearing as a useful descriptive fiction layered over more fundamental realities. If relation can sometimes belong to the order of what first organizes appearance, then harmonic structure becomes a plausible candidate for primitive description rather than an ornamental afterthought.

From the side of signification, a parallel displacement appears in another idiom. Derrida's account of différance interrupts the fantasy of self-present meaning by insisting that what matters in signification is never fully gathered into a term's supposed identity. In the essay by that name, différance is described as what "is never offered to the present" \parencite{derrida_1982}. Sense therefore depends on spacing, deferral, trace, and differential articulation rather than on pure presence or origin. Read in this register, relation is not an accident that later befalls meaning. It belongs to the conditions under which meaning becomes possible. The sign bears the mark of what it is not; it receives force from intervals, delays, and contrasts that prevent it from coinciding transparently with itself.

Lacan radicalizes the point in a way that is especially pertinent here. His formula that a signifier represents a subject for another signifier refuses the idea that the subject could be located intact at either end of an already completed exchange \parencite{lacan_1966}. The subject is not reducible to \texttt{S1} or \texttt{S2}; it appears as a barred effect of the chain, and of the chain's insertion in the field of the Other. Communication, on this view, is not simply the transport of content between fully present interiors. It is a differential process in which position, address, and intelligibility emerge through articulation. That does not convert HIT into psychoanalysis, nor should it. It does, however, sharpen one of the chapter's central points: informational significance need not be thought as a property lying dormant inside isolated units. It may arise at the level of relation, interval, and sequence.

Jung names a comparable organizing tension within psychology, though in a different vocabulary. The Self, in the analytic tradition, is neither identical with the ego nor simply another content hidden somewhere in the psyche. It names an integrating principle that becomes legible through the dynamic relation between conscious and unconscious life rather than through the sovereignty of either pole taken alone \parencite{jung_1959}. Von Franz makes the point more operational when she describes the Self as an orienting factor irreducible to the conscious personality and visible only through processes of individuation rather than through static introspection \parencite{vonfranz_1964}. Read cautiously, this does not import clinical doctrine into ontology. It shows that even in depth psychology, what organizes the field may be neither a fixed substance nor a self-transparent subject, but a process of relation, tension, and gradual integration.

Lotman's semiosphere gives this relationality a broader semiotic topology. No isolated language or code stands first and only later enters relation; the semiosphere functionally precedes its local formations, and meaning takes shape through boundaries, translations, filters, and asymmetric passages between heterogeneous spaces \parencite{lotman_2005}. Peirce and Bateson deepen the same orientation from adjacent directions. If the sign is irreducibly triadic, and if information appears where a difference makes a difference, then relation cannot be treated as residue. It belongs to the conditions under which distinction, transmission, and effect become thinkable at all \parencite{peirce_1931,bateson_1972}. This does not imply that every relation is informative or that every pattern is semiotic. It implies something more precise: some structured relations may be prior to the units later abstracted from them, and the explanatory burden must sometimes begin there.

Nietzsche's "On Truth and Lies in a Nonmoral Sense" introduces a necessary counterpressure into the entire discussion. Concepts appear there as worn metaphors, and truths as stabilized conventions that have forgotten their figurative genesis rather than as transparent mirrors of self-identical essences \parencite{nietzsche_1873}. The warning does not collapse ontology into arbitrariness. It blocks the temptation to mistake classificatory names for the last word on what there is. Wheeler reopens the same question from another side when he proposes that information may not be mere bookkeeping applied to an already complete world, but part of what makes such a world intelligible in the first place \parencite{wheeler_1990}. Read together, Nietzsche and Wheeler mark a narrow passage: one can take information seriously without reifying every formal structure into metaphysical finality.

The ontological wager of this chapter can therefore be stated more carefully. If relation can be primary in existential, semiotic, and relational-physical vocabularies, then it is not conceptually illicit to ask whether harmonic relations among oscillatory or mode-bearing processes might count as primitive informational structures. The task is neither to reduce the world to one acoustic image nor to inflate formal analogy into exhaustive identity. It is to test whether interval, ratio, recurrence, and patterned coupling belong often enough to the order of first intelligibility that a theory of harmonic information becomes worth building. At that point the burden of proof is no longer conceptual permission. It is empirical reach and theoretical precision.


\section{The interval as a minimal informational carrier}


Within that relational ontology, the interval becomes the candidate unit of interest. An interval, in the sense relevant to HIT, is not primarily a named musical object such as a fifth or a third. It is a ratio relation between oscillatory frequencies or, more generally, between modes capable of entering patterned coupling. In acoustics that relation becomes audible as interval. In other domains it may appear as a ratio between frequencies, periodicities, eigenmodes, or recurrent scales of organization. Given two frequencies, \texttt{f1} and \texttt{f2}, the interval may be written as \texttt{r = f1:f2}. What matters about this relation is not only that two magnitudes stand in comparison, but that their coupling generates an interference pattern with distinctive structural properties.

The first of those properties is scale robustness. If both frequencies are multiplied by the same constant \texttt{k}, the absolute frequencies change but the ratio does not. More carefully stated, the interval's relational organization is preserved up to temporal rescaling. This does not mean that every physical system will respond identically across all scales. Media, thresholds, and receptor dynamics still matter. It means something weaker and crucial: the interval is not defined by absolute pitch, but by proportional relation. That alone makes it a stronger candidate for cross-domain description than any frequency taken in isolation.

The second property is pattern generation. Intervals are not abstract ratios floating above the material world; they generate concrete interference organizations. In visual form, Lissajous figures make this especially intuitive. When perpendicular oscillations stand in simple relations, the resulting curves exhibit repeatable closure, symmetry, and topological economy; when the relations become more complex, the curves become correspondingly less simple and less recurrent \parencite{lissajous_1857,gallozzi_2023}. In driven media, related phenomena appear in standing-wave and Faraday-like pattern formation, where periodic forcing organizes matter into stable spatial regimes \parencite{engels_2007,nguyen_2019}. The visual case does not prove the ontological thesis, but it reveals something important: relation can become visible as form. Cymatic traditions, though often rhetorically overstated, are useful at minimum as reminders that oscillatory relations can leave stable traces in matter and geometry \parencite{jenny_1967}.

The third property is dynamical privilege. In models of coupled oscillation, some ratios are more likely than others to anchor stable mode-locking or recurrence structure, especially once one looks at Arnold tongues, locking regimes, and the broader synchronization literature rather than at isolated equations alone \parencite{kuramoto_1984,strogatz_2000}. Reviews such as Acebron and colleagues help situate this broader terrain of synchronization and partial locking without themselves carrying the entire argument about ratio privilege \parencite{acebron_2005}. In neural and cortical systems, ratios likewise appear not merely as incidental numerical coincidences but as organizing constraints within coupling architectures \parencite{canolty_2010,jacobs_2025}. This does not license the naive claim that "simple ratios are always better." It does support a more exact claim: certain intervallic relations are repeatedly recruited by systems because they afford recurrent and stable coordination. That is exactly the sort of condition under which an informational role becomes plausible.

The fourth property is differential interpretability. An interval functions as an informational carrier only if a system can distinguish it from alternative relations in a stable way. This point returns us to Bateson. A ratio is not informative simply because it can be written. It becomes informative when it constrains behavior, coupling, prediction, or response differently than another ratio would. In this sense, the interval is a candidate minimal carrier because it is the smallest relational unit capable of generating a non-trivial difference in oscillatory organization. A single oscillator can oscillate. Two oscillators can enter into relation. From relation comes interference; from interference, pattern; from pattern, the possibility of differential uptake.

At this point, the central proposal can be formulated more precisely. HIT treats the interval as a minimal informational carrier not because it is musically familiar, but because it satisfies four conditions at once: it is relational rather than substantial, robust under common scaling, generative of recurrent structure, and in principle capable of differential interpretation by receiving systems. None of these conditions alone is sufficient. Together they make the interval the strongest initial candidate for the ontological primitive that this program requires.


\section{A minimal ontological proposition for HIT}


The argument of this chapter can now be compressed into a minimal proposition. Information, in the sense relevant to HIT, should be understood as structured relation with the capacity to generate repeatable difference for a system. Harmonic information is the subset of such information that arises through patterned relations among oscillatory, wave-bearing, or mode-organized processes. An interval is the minimal candidate carrier of that information because it specifies a ratio relation whose organization is preserved beyond any single absolute instantiation and whose interference pattern can, under suitable conditions, be stabilized, distinguished, and acted upon.

This formulation remains deliberately narrower than the broader ambitions often attached to harmonic discourse. It does not yet imply that consonance is universally valued, that harmonic simplicity is always biologically optimal, or that every domain we organize as reality is governed by the same ratio laws. Those are later questions. The present chapter has only secured the weaker but necessary claim that a relational ontology of information is coherent, that oscillatory relations offer a plausible site for such an ontology, and that the interval can be treated as the minimal conceptual unit around which the rest of the framework may be organized.

It is useful, finally, to separate three levels of commitment. The observation is that multiple literatures repeatedly assign explanatory weight to structured harmonic relations. The hypothesis is that intervals can function as minimal informational carriers across heterogeneous systems. The inference, more cautious than a doctrine but stronger than a metaphor, is that HIT should be built on a relational ontology rather than on a substance-first model of information. Everything that follows depends on preserving that distinction.

If this is correct, then consonance cannot be defined as the intrinsic property of a ratio taken in abstraction from medium, timbre, body, and context. Once information is treated as relational, consonance must be treated relationally as well. That is the task of the next chapter.
